\section{Why the uniformizing lattices have equal covolume}
\label{sec:selfdual-covol}
Here we give the index argument used in \eqref{eq:volume-flag}.
Let \(E/F\) be the CM extension with \(F=E^+\) and involution
\(x\mapsto\bar x\). Let \(r=p+q\), with \(p,q\ge1\).
Write \(G\) for the matrix of a nondegenerate Hermitian form.
At the deep level the source takes the positive integer \(M=l.M'\)
(the combined congruence modulus), with \(M\ge3\).
At the chosen embedding, the components under comparison have the
same CM type, signature and archimedean measure.

\paragraph{Reduction to two lattices on one Hermitian space.}
The full-level presentation of each component preserves its covolume
and supplies a trace-self-dual lattice. To pass from two presentations
to a single ambient space the source uses a rational similitude
\(P\) supplied by the classification of the polarized data. In its
normalization, for a move from \(X\) to \(Y\) the ratio of Hermitian
forms satisfies
\[
 h_Y(Pv,Pw)=\frac{\kappa_X}{\kappa_Y}\,h_X(v,w).
\]
Thus the scalar-weighted trace pairing is carried exactly into the
second one, the lattice becomes \(P\Lambda\), and reduction modulo
\(M(P\Lambda)\) is transported by \(P\). In particular
\(P\Gamma P^{-1}\) is the transported group. Its flag-domain
identification preserves the invariant volume. In what follows the
two full lattices, say \(\Lambda,\Lambda'\), have the same \(G\),
CM type, trace parameter and signature. The local self-dual isometry
theorem, assembled over the finite places, supplies
\(g\in U(G)(\mathbb A_{F,f})\) with
\(g\widehat\Lambda=\widehat{\Lambda'}\). The proof only needs a
finite-adelic isometry at this point.

\paragraph{Congruence groups and intersections.}
Write \(\iota\) for the diagonal embedding of rational matrices into
finite-adelic matrices. Let \(K\) be the subgroup of determinant-one
isometries preserving \(\widehat\Lambda\) which act trivially modulo
\(M\widehat\Lambda\). Put \(K'=gKg^{-1}\). The full-level congruence
description identifies \(\Gamma=\iota^{-1}K\) and
\(\Gamma'=\iota^{-1}K'\) with the two uniformizing groups. Here
isometries are taken in the same ambient rational unitary group.
In particular this description applies the determinant-one condition
to the level group; \(g\) can have arbitrary norm-one determinant.

We record how intersection indices pass to the rational points.
If \(J=K\cap K'\), reduction to cosets gives injective maps
\[
 \Gamma/(\Gamma\cap\Gamma')\longrightarrow K/J,\qquad
 \Gamma'/(\Gamma\cap\Gamma')\longrightarrow K'/J ,
\]
where each quotient denotes cosets and carries no normality assumption.
Both maps are surjective. Indeed \(J\) contains a deeper principal
congruence subgroup. Explicitly, denominators of \(g,g^{-1}\)
are bounded on the completed lattice by positive integers \(b,a\).
Then any isometry congruent to the identity modulo \(Mab\) on
\(\widehat\Lambda\) lies in \(K\) and its conjugate back by \(g\)
is congruent to the identity modulo \(M\). This follows by applying
\(g^{-1}(s-1)g\) to a lattice vector and cancelling the two denominators;
the condition for the inverse follows from the congruence subgroup
property. The source obtains rational isotropy of \(G\) from indefinite
signature and the trace-self-dual data and applies strong approximation
for its special unitary group. It gives, for every \(s\in K\cup K'\),
a rational isometry \(\gamma\) with
\(\iota(\gamma)^{-1}s\) in this deeper subgroup. Thus each of the
displayed cosets has a rational representative. The same argument works
with \(g\) replaced by a power, or by a rational isometry.

\paragraph{The ratio of two indices.}
For commensurable groups \(A,B\) write
\[
 r(A,B)=\frac{[A:A\cap B]}{[B:A\cap B]} .
\]
If a third group is commensurable with both, one can compute all ratios
using their triple intersection. This gives \(r(A,B)r(B,C)=r(A,C)\).
Conjugation carries a pair of indices into a pair of the same sizes.
Consequently \(r(K,g^jKg^{-j})=r(K,K')^j\) for positive \(j\).
The indices here are finite since the groups intersect in a deeper open
congruence subgroup of the compact level group.

Put \(z=\overline{\det g}\) and \(s=(zI)g^r\).
Unitarity gives \(z\det g=1\), so \(zI\) is unitary and
\(\det s=z^r(\det g)^r=1\). The central scalar drops out under
conjugation. Apply strong approximation at level \(M\) to \(s\), writing
\(s=\iota(\gamma)k\) for a rational unitary \(\gamma\) and \(k\in K\).
Then
\[
 g^rKg^{-r}=sKs^{-1}=\iota(\gamma)K\iota(\gamma)^{-1}.
\]
The corresponding rational groups are \(\Gamma\) and
\(\gamma\Gamma\gamma^{-1}\). Their quotient volumes in the archimedean
flag domain are equal by the isometry induced by \(\gamma\).
For two commensurable uniformizing groups \(A,B\) acting freely at
this level, the common covering for \(A\cap B\) has volume both
\([A:A\cap B]\operatorname{covol}(A)\) and
\([B:A\cap B]\operatorname{covol}(B)\), with identical measure
normalization. Applying this to the two rationally conjugate groups,
and using the coset bijections, gives \(r(K,g^rKg^{-r})=1\).
The ratio \(r(K,K')\) is positive, so its \(r\)-th power being one
implies \(r(K,K')=1\). Using the coset bijections for \(g\) and then
the common-cover formula yields
\(\operatorname{covol}(\Gamma)=\operatorname{covol}(\Gamma')\).

\paragraph{Source and scope.}
The two reductions to full-level data and a common Hermitian form are
\srcid{ExistsFullLevelGroupDataSameCovolume} and
\srcid{ExistsTransportSelfDualDataOneHermitianSpace}.
The index proof is \srcid{CovolumeEqSelfDualLatticesSameHermitianSpace};
the component application is \srcid{ComplexUniformisationComponentsCovolumeEq}.
The input theorems here are the full-level uniformization, rational
similitude and isotropy, finite-local equivalence of the relevant
trace-self-dual lattices, and special unitary strong approximation.
